\documentclass[12pt,a4paper]{article}

% --- PAQUETES DE CONFIGURACIÓN ---
\usepackage[spanish]{babel}
\usepackage[utf8]{inputenc}
\usepackage[T1]{fontenc}
\usepackage{geometry}
\geometry{margin=2.5cm}
\usepackage{listings}
\usepackage{xcolor}
\usepackage{hyperref}
\usepackage{endnotes}
\usepackage{amsmath}

% --- CONFIGURACIÓN DE APARIENCIA PARA CÓDIGO ---
\definecolor{codegreen}{rgb}{0,0.6,0}
\definecolor{codegray}{rgb}{0.5,0.5,0.5}
\definecolor{codeblue}{rgb}{0,0,0.8}
\definecolor{backcolour}{rgb}{0.96,0.96,0.96}

\lstdefinestyle{mystyle}{
    backgroundcolor=\color{backcolour},
    commentstyle=\color{codegreen},
    keywordstyle=\color{codeblue},
    numberstyle=\tiny\color{codegray},
    stringstyle=\color{red},
    basicstyle=\ttfamily\small,
    breakatwhitespace=false,
    breaklines=true,
    captionpos=b,
    keepspaces=true,
    numbers=left,
    numbersep=5pt,
    showspaces=false,
    showstringspaces=false,
    showtabs=false,
    tabsize=2,
    literate=
        {á}{{\'a}}1
        {é}{{\'e}}1
        {í}{{\'i}}1
        {ó}{{\'o}}1
        {ú}{{\'u}}1
        {ñ}{{\~n}}1
}
\lstset{style=mystyle}

\lstdefinelanguage{json}{language=Java}
\lstdefinelanguage{properties}{
  sensitive=false,
  morecomment=[l]{\#},
  morecomment=[l]{!}
}

% --- DATOS DEL DOCUMENTO ---
\title{Bases de Datos Relacionales de Alto Rendimiento: Amazon Aurora}
\date{\normalsize\today}

\begin{document}

\maketitle

\begin{abstract}

\end{abstract}

{
\let\clearpage\relax
\tableofcontents
}
\newpage

\section{Introducción}

Amazon Aurora es un motor de base de datos relacional compatible con MySQL y PostgreSQL diseñado específicamente para la nube, que ofrece hasta 5 veces el rendimiento de MySQL estándar y 3 veces el de PostgreSQL.

Redefine la arquitectura de las bases de datos relacionales tradicionales al separar por completo la capa de cómputo (\textit{Database Nodes}) de la capa de almacenamiento (\textit{Distributed Storage Cluster}).

En Event Hub el backend Spring Boot (\texttt{eventhub-back-springboot}) deja de usar H2 en memoria y persiste eventos y compras en un clúster Aurora PostgreSQL. El despliegue se automatiza con \texttt{make deploy}: primero \texttt{ec2.sh} (instancia y grupo de seguridad) y a continuación \texttt{aurora.sh} (clúster, instancia de cómputo y secreto de usuario \textit{master}). No hay un script aparte para Secrets Manager: la contraseña la crea Aurora al crear el clúster.

\subsection{Almacenamiento Desacoplado y Tolerancia a Fallos}

A diferencia de los motores relacionales convencionales en los que cada nodo gestiona su propio volumen de disco, Aurora utiliza un volumen de almacenamiento virtualizado y auto-escalable:

\begin{itemize}
    \item \textbf{Replicación en 3 Zonas de Disponibilidad (AZ)}: Los datos se dividen en segmentos de 10 GB y se replican 6 veces a lo largo de 3 Zonas de Disponibilidad (2 copias por AZ).
    \item \textbf{Modelo de Quórum de Escritura/Lectura}: Para procesar una escritura se requiere un quórum de 4 de 6 copias ($4/6$), garantizando la disponibilidad de lectura ante la pérdida de una AZ completa más una copia adicional ($3/6$).
    \item \textbf{Auto-Reparación (\textit{Self-Healing})}: La capa de almacenamiento escanea de forma continua los bloques de disco en busca de errores y repara automáticamente los segmentos dañados en segundo plano sin impacto en la latencia.
    \item \textbf{Crecimiento Automático}: El almacenamiento se expande dinámicamente según la demanda hasta un límite de 128 TiB por clúster, sin necesidad de aprovisionar capacidad por adelantado.
\end{itemize}

\subsection{Aurora Serverless v2}

Aurora Serverless v2 permite escalar automáticamente los recursos de cómputo (vCPU y memoria RAM) en fracciones de segundo. La capacidad se mide en Unidades de Capacidad de Aurora (\textit{Aurora Capacity Units} o ACUs)%
\endnote{Una ACU representa aproximadamente 2 GiB de memoria RAM junto con sus correspondientes recursos de vCPU y red.}, lo que permite ajustar dinámicamente la infraestructura ante picos impredecibles de tráfico sin afectar las conexiones activas.

El laboratorio usa una instancia aprovisionada \texttt{db.t3.medium}, clase mínima habitual de Aurora PostgreSQL en AWS Academy. Serverless v2 no se emplea en los scripts, pero el mismo clúster podría migrarse a \texttt{db.serverless} sin cambiar el modelo de almacenamiento.

\section{Alta Disponibilidad, Escalabilidad de Lectura y Seguridad}

El diseño de Amazon Aurora optimiza los tiempos de recuperación y minimiza el retraso de replicación en entornos de producción.

\subsection{Réplicas de Lectura y Recuperación de Desastres}
\begin{itemize}
    \item \textbf{Hasta 15 Réplicas de Lectura}: Aurora permite conectar hasta 15 réplicas que comparten el mismo volumen de almacenamiento subyacente, reduciendo el retardo de replicación a milisegundos.
    \item \textbf{Failover Automático}: Si el nodo primario de escritura falla, Aurora promueve automáticamente una de las réplicas de lectura a nodo primario en menos de 30 segundos.
    \item \textbf{Global Databases}: Permite la replicación entre diferentes regiones de AWS con latencias menores a 1 segundo, ofreciendo recuperación ante desastres a nivel global.
\end{itemize}

\subsection{Estrategia de Cifrado y Aislamiento de Red}
El clúster de Aurora de Event Hub no es accesible desde Internet (\texttt{--no-publicly-accessible}). Comparte el grupo de seguridad de la instancia EC2 y solo admite PostgreSQL desde ese mismo grupo:
\begin{itemize}
    \item \textbf{Cifrado en Reposo}: Utiliza algoritmos AES-256 gestionados a través de claves de \textit{AWS KMS}, asegurando datos, registros de transacciones (\textit{WAL/binlogs}), copias de seguridad e instantáneas (\textit{snapshots}).
    \item \textbf{Cifrado en Tránsito}: Fuerza el uso de TLS/SSL en las comunicaciones cliente-servidor.
    \item \textbf{Grupo de seguridad compartido}: \texttt{ec2.sh} crea \texttt{eventhub-sg} y escribe su \texttt{GroupId} en \texttt{aws-scripts/sg-id.txt}. \texttt{aurora.sh} lee ese fichero y abre el puerto 5432 con origen el propio grupo, de modo que la EC2 puede conectar al clúster sin un segundo \textit{security group}.
\end{itemize}

\section{Gestión Centralizada de Secretos con AWS Secrets Manager}

Para evitar la exposición de credenciales maestras en el código, Event Hub no genera la contraseña en el portátil ni la pasa a RDS por la CLI. Al crear el clúster, Aurora genera la password y la guarda en Secrets Manager (\texttt{--manage-master-user-password}).

El JSON del secreto gestionado usa las mismas claves que Spring Boot espera (\texttt{username}, \texttt{password}, \texttt{host}, \texttt{port} y, si RDS la incluye, \texttt{dbname}). El \texttt{host} solo es válido cuando el clúster ya tiene \textit{endpoint}: lo rellena Aurora, no el script.

RDS no permite elegir el nombre del secreto. Lo asigna AWS con el prefijo \texttt{rds!cluster-}\ldots{} Por eso \texttt{aurora.sh} no puede dejar un nombre fijo en \texttt{application.properties}. Cuando el clúster está \texttt{available}, consulta el ARN del secreto, obtiene su \texttt{Name} con \texttt{describe-secret} (esa llamada \emph{no} devuelve la password) y escribe \texttt{aurora.properties}, el mismo patrón que \texttt{cognito.properties}.

No se usa \texttt{get-secret-value} ni \texttt{--master-user-password}: si el CLI leyera la password y la reenviara a RDS, viajaría por la red hasta el portátil. Con \texttt{--manage-master-user-password} la clave no sale de AWS.

La rotación del secreto \textit{master} la gestiona Aurora (no una Lambda del laboratorio). Puede forzarse con \texttt{aws rds modify-db-cluster --rotate-master-user-password --apply-immediately}. El detalle de rotación con dos usuarios de aplicación queda fuera de este documento. Al borrar el clúster, RDS elimina también el secreto gestionado.

\section{Despliegue CLI: Clúster de Amazon Aurora PostgreSQL}

La secuencia está en \texttt{aws-scripts/aurora.sh}. \texttt{make deploy} lo ejecuta después de \texttt{ec2.sh}. Cada comando se explica a continuación.

\subsection{Paso 1: Subnet group en la VPC por defecto}

Aurora exige subnets en al menos dos zonas de disponibilidad. El laboratorio toma las subnets de la VPC por defecto de AWS Academy (este es el primer paso de red en \texttt{aurora.sh create}):

\begin{lstlisting}[language=bash,basicstyle=\ttfamily\footnotesize]
# --filters: VPC marcada como default
# --query: extrae el VpcId en texto plano
VPC_ID=$(aws ec2 describe-vpcs \
    --filters Name=isDefault,Values=true \
    --query 'Vpcs[0].VpcId' \
    --output text)

SUBNET_IDS=$(aws ec2 describe-subnets \
    --filters "Name=vpc-id,Values=$VPC_ID" \
    --query 'Subnets[].SubnetId' \
    --output text)

# --db-subnet-group-name: nombre del grupo (Aurora exige >= 2 AZ)
# --subnet-ids: subnets de la VPC por defecto
aws rds create-db-subnet-group \
    --db-subnet-group-name eventhub-aurora-subnets \
    --db-subnet-group-description "Subnets Aurora EventHub" \
    --subnet-ids $SUBNET_IDS
\end{lstlisting}

\vspace{1em}

\texttt{--query} usa JMESPath: \texttt{Vpcs[0].VpcId} es el primer (y único) resultado; \texttt{Subnets[].SubnetId} es la lista de identificadores. \texttt{\$SUBNET\_IDS} no va entrecomillado a propósito: la CLI espera varios argumentos \texttt{subnet-xxxx}, no una sola cadena con espacios.

El identificador del clúster lo elige \texttt{aurora.sh} (\texttt{eventhub-cluster-aurora}) y lo escribe en \texttt{db-cluster-id.txt}. Ese fichero solo lo usa \texttt{aurora.sh} (por ejemplo al borrar).

Constantes que el script sí conoce de antemano (no salen de un secreto previo):

\begin{lstlisting}[language=bash,basicstyle=\ttfamily\footnotesize]
DB_CLUSTER_ID="eventhub-cluster-aurora"
DB_USER="master"
DB_NAME="eventhub"
DB_ENGINE="aurora-postgresql"
DB_PORT=5432
\end{lstlisting}

\vspace{1em}

\begin{itemize}
    \item \texttt{DB\_CLUSTER\_ID}: nombre del clúster. Lo elige el laboratorio, no AWS.
    \item \texttt{DB\_USER}: usuario \textit{master}. En Aurora PostgreSQL no puede ser \texttt{admin} (palabra reservada del motor); en esta versión el nombre elegido es \texttt{master}.
    \item \texttt{DB\_NAME}: base inicial (\texttt{--database-name}). Coincide con el \texttt{dbname} de la URL JDBC cuando el secreto lo incluye.
    \item \texttt{DB\_ENGINE}: valor \texttt{aurora-postgresql} que RDS espera en \texttt{create-db-cluster --engine}.
    \item \texttt{DB\_PORT}: 5432, puerto por defecto de PostgreSQL.
\end{itemize}

\subsection{Paso 2: Transmitir el grupo de seguridad desde EC2}

Aurora no crea un \textit{security group} propio. \texttt{ec2.sh} ya ha creado \texttt{eventhub-sg} (HTTP, HTTPS y SSH) y escribe el \texttt{GroupId} en un fichero, el mismo patrón que \texttt{public-ip.txt} o \texttt{callback-url.txt}:

\begin{lstlisting}[language=bash,basicstyle=\ttfamily\footnotesize]
# GroupId del SG compartido con Aurora
echo "$SG_ID" > "$SG_ID_FILE"
\end{lstlisting}

\vspace{1em}

\texttt{\$SG\_ID} es el identificador que devolvió \texttt{aws ec2 create-security-group} (por ejemplo \texttt{sg-0abc123}). \texttt{\$SG\_ID\_FILE} apunta a \texttt{aws-scripts/sg-id.txt}. El fichero está en \texttt{.gitignore} (\texttt{aws-scripts/*.txt}): no se sube al repositorio.

\texttt{aurora.sh} lo lee así:

\begin{lstlisting}[language=bash,basicstyle=\ttfamily\footnotesize]
SG_ID=$(< "$SG_ID_FILE")
SG_ID="${SG_ID%%$'\n'}"
\end{lstlisting}

\vspace{1em}

\texttt{\$(< fichero)} es la redirección de Bash que carga el contenido del fichero en una variable, sin lanzar \texttt{cat}. La segunda línea elimina un posible salto de línea final, para que \texttt{--group-id} reciba solo el identificador.

Con ese \texttt{GroupId} se abre PostgreSQL. El tráfico permitido no es ``cualquier IP'', sino el propio grupo: la instancia EC2 y el clúster Aurora están en el mismo SG, y el origen \texttt{--source-group} autoriza a los miembros de ese grupo.

\begin{lstlisting}[language=bash,basicstyle=\ttfamily\footnotesize]
# --group-id: SG compartido con EC2 (leido de sg-id.txt)
# --protocol / --port: PostgreSQL
# --source-group: la instancia EC2 del mismo SG puede conectar a Aurora
aws ec2 authorize-security-group-ingress \
    --group-id "$SG_ID" \
    --protocol tcp \
    --port 5432 \
    --source-group "$SG_ID"
\end{lstlisting}

\vspace{1em}

Si la regla 5432 ya existe (un segundo \texttt{aurora.sh create}), el script no vuelve a crearla.

\subsection{Paso 3: Crear el clúster con password gestionada por Aurora}

\begin{lstlisting}[language=bash,basicstyle=\ttfamily\footnotesize]
# --db-cluster-identifier: nombre elegido por aurora.sh
# --engine: aurora-postgresql
# --master-username: master (no admin)
# --manage-master-user-password: Aurora crea el secreto; la password no viaja por la CLI
# --database-name: base inicial
# --vpc-security-group-ids: mismo SG que la instancia EC2
# --db-subnet-group-name: subnets en al menos 2 AZ
# --backup-retention-period 1: retencion minima (laboratorio)
# --no-deletion-protection: permite borrar el cluster con make delete
aws rds create-db-cluster \
    --db-cluster-identifier eventhub-cluster-aurora \
    --engine aurora-postgresql \
    --master-username master \
    --manage-master-user-password \
    --database-name eventhub \
    --port 5432 \
    --vpc-security-group-ids "$SG_ID" \
    --db-subnet-group-name eventhub-aurora-subnets \
    --backup-retention-period 1 \
    --no-deletion-protection
\end{lstlisting}

\vspace{1em}

\texttt{--manage-master-user-password} es incompatible con \texttt{--master-user-password}: o bien RDS gestiona la clave, o bien el CLI la envía. El laboratorio elige lo primero. RDS crea un secreto cuyo nombre tiene el prefijo \texttt{rds!cluster-} y rellena \texttt{username}, \texttt{password}, \texttt{engine}, \texttt{port} y, cuando el clúster tiene DNS, \texttt{host}.

\texttt{aws rds wait db-cluster-available} bloquea hasta que el estado del clúster es \texttt{available}. Sin ese wait, \texttt{create-db-instance} fallaría porque el clúster aún no está listo. Puede tardar varios minutos.

\subsection{Paso 4: Instancia de cómputo}

Un clúster Aurora sin instancia de cómputo no acepta conexiones JDBC. Hay que crear al menos un nodo:

\begin{lstlisting}[language=bash,basicstyle=\ttfamily\footnotesize]
# --db-instance-identifier: nodo de computo del cluster
# --db-cluster-identifier: cluster del paso anterior
# --db-instance-class: clase minima habitual de Aurora PostgreSQL
# --no-publicly-accessible: solo accesible desde la VPC (la EC2)
aws rds create-db-instance \
    --db-instance-identifier eventhub-aurora-instance \
    --db-cluster-identifier eventhub-cluster-aurora \
    --engine aurora-postgresql \
    --db-instance-class db.t3.medium \
    --no-publicly-accessible

aws rds wait db-instance-available \
    --db-instance-identifier eventhub-aurora-instance
\end{lstlisting}

\vspace{1em}

\texttt{db.t3.medium} es la clase que admite AWS Academy para Aurora PostgreSQL. \texttt{--no-publicly-accessible} impide un DNS público: Spring Boot, en la EC2, llega al clúster por la red privada de la VPC.

\subsection{Paso 5: Nombre del secreto y \texttt{aurora.properties}}

Cuando la instancia está \texttt{available}, el clúster ya tiene \textit{endpoint} de escritura y el secreto gestionado incluye el \texttt{host}. El script no lee el \texttt{SecretString}: solo necesita el nombre para que Spring Boot importe ese secreto.

\begin{lstlisting}[language=bash,basicstyle=\ttfamily\footnotesize]
# --query: ARN del secreto que creo Aurora (MasterUserSecret)
# describe-secret no devuelve la password
SECRET_ARN=$(aws rds describe-db-clusters \
    --db-cluster-identifier eventhub-cluster-aurora \
    --query 'DBClusters[0].MasterUserSecret.SecretArn' \
    --output text)

SECRET_NAME=$(aws secretsmanager describe-secret \
    --secret-id "$SECRET_ARN" \
    --query Name \
    --output text)
\end{lstlisting}

\vspace{1em}

El \textbf{ARN} (\textit{Amazon Resource Name}) es el identificador único del recurso en AWS: servicio, región, cuenta y nombre. Ejemplo: \texttt{arn:aws:secretsmanager:us-east-1:123456789012:secret:rds!cluster-\ldots}. El \texttt{Name} es la parte que Spring Cloud AWS usa en \texttt{aws-secretsmanager:}.

Ese nombre no es conocido en tiempo de compilación. \texttt{aurora.sh} genera un properties, igual que \texttt{cognito.sh} genera el emisor JWT:

\begin{lstlisting}[language=properties,basicstyle=\ttfamily\footnotesize]
# Generado por aws-scripts/aurora.sh. No editar a mano: se sobrescribe en cada create.
# El secreto lo crea Aurora; el host solo es válido cuando el clúster ya existe.
spring.config.import=aws-secretsmanager:rds!cluster-xxxxxxxxxxxx
\end{lstlisting}

\vspace{1em}

El fichero se escribe en \texttt{eventhub-back-springboot/src/main/resources/aurora.properties} y está en el \texttt{.gitignore} de ese módulo. Al hacer \texttt{aurora.sh delete} se elimina junto con \texttt{db-cluster-id.txt}.

\section{Integración con Spring Boot}

El backend declara el driver PostgreSQL (Aurora PostgreSQL habla el protocolo de PostgreSQL; no hay un driver ``Aurora'') y el \textit{starter} de Spring Cloud AWS para Secrets Manager. H2 queda restringido a los tests (\texttt{scope=test}), de modo que \texttt{mvn test} no necesita Aurora.

\begin{lstlisting}[language=XML,basicstyle=\ttfamily\footnotesize]
<dependency>
    <groupId>org.postgresql</groupId>
    <artifactId>postgresql</artifactId>
    <scope>runtime</scope>
</dependency>
<dependency>
    <groupId>io.awspring.cloud</groupId>
    <artifactId>spring-cloud-aws-starter-secrets-manager</artifactId>
</dependency>
\end{lstlisting}

\vspace{1em}

La versión del \textit{starter} la fija el BOM \texttt{spring-cloud-aws-dependencies} (4.1.0), compatible con Spring Boot 4.1. El \textit{starter} registra un \texttt{ConfigDataLocationResolver} que reconoce el prefijo \texttt{aws-secretsmanager:}.

Las claves de nivel superior del JSON del secreto se convierten en propiedades del entorno de Spring. No es necesario \texttt{spring.cloud.aws.region.static}: el SDK usa \texttt{DefaultAwsRegionProviderChain} (\texttt{AWS\_REGION}, \texttt{\~/.aws/config} o metadatos de EC2).

\texttt{application.properties} no nombra el secreto de RDS. Importa el fichero generado:

\begin{lstlisting}[language=properties,basicstyle=\ttfamily\footnotesize]
spring.config.import=classpath:cognito.properties,classpath:aurora.properties

spring.datasource.username=${username}
spring.datasource.password=${password}
spring.datasource.url=jdbc:postgresql://${host}:${port}/${dbname:eventhub}

spring.jpa.hibernate.ddl-auto=update
\end{lstlisting}

\vspace{1em}

\begin{itemize}
    \item \texttt{spring.config.import}: \texttt{classpath:cognito.properties} aporta el emisor JWT; \texttt{classpath:aurora.properties} aporta, a su vez, \texttt{aws-secretsmanager:} con el nombre \texttt{rds!cluster-\ldots} que asignó RDS.
    \item \texttt{spring.datasource.username=\$\{username\}}: mapea la clave AWS \texttt{username} a la propiedad estándar de Spring Boot. Igual para \texttt{password}.
    \item \texttt{spring.datasource.url}: se construye con \texttt{host}, \texttt{port} y \texttt{dbname}. El esquema es \texttt{jdbc:postgresql://} (el driver de PostgreSQL), no \texttt{jdbc:\$\{engine\}://}: en el secreto gestionado \texttt{engine} suele valer \texttt{postgres}, y \texttt{aurora-postgresql} tampoco es un esquema JDBC válido.
    \item \texttt{\$\{dbname:eventhub\}}: si el JSON de RDS no incluye \texttt{dbname}, Spring usa \texttt{eventhub}, la base creada con \texttt{--database-name}.
    \item \texttt{spring.jpa.hibernate.ddl-auto=update}: Hibernate crea o altera tablas (\texttt{evento}, \texttt{compra}) al arrancar.%
\endnote{En el laboratorio la aplicación conecta como usuario \textit{master} (\texttt{master}), lo que permite ese DDL. En un entorno profesional el esquema se gestionaría con migraciones (por ejemplo Flyway o Liquibase) y un usuario de aplicación sin privilegios de administración.}
\end{itemize}

El recurso que ejecuta Spring Boot (la EC2) debe poder invocar \texttt{secretsmanager:GetSecretValue} sobre el secreto \texttt{rds!cluster-\ldots}. En el laboratorio eso lo cubren las credenciales de AWS Academy o un perfil de instancia; el detalle de IAM está en la guía correspondiente.

\section{Conclusión}

Amazon Aurora ofrece una arquitectura relacional de clase empresarial: almacenamiento distribuido con quórum, cómputo independiente del disco y failover en segundos. En Event Hub esa arquitectura se concreta en un clúster Aurora PostgreSQL cuya contraseña genera el propio RDS al crear el clúster (\texttt{--manage-master-user-password}), de modo que la clave no transita por la CLI. El nombre del secreto lo asigna AWS; \texttt{aurora.sh} lo deja en \texttt{aurora.properties} para que Spring Boot monte el \textit{datasource} sin credenciales en el repositorio. El grupo de seguridad de la EC2, transmitido por \texttt{sg-id.txt}, aisla el puerto 5432 al interior de la VPC. El resultado es un entorno alineado con las prácticas de secretos, red y persistencia exigibles en una aplicación desplegada en AWS.

\bibliographystyle{plain}
\theendnotes

\end{document}
